\documentclass[11pt]{article}
\usepackage[a4paper,margin=27mm]{geometry}
\usepackage{amsmath,amssymb,amsthm}
\usepackage[T1]{fontenc}
\usepackage{lmodern}
\usepackage{microtype}
\usepackage[hidelinks]{hyperref}
\newtheorem{lemma}{Lemma}
\newtheorem{theorem}{Theorem}
\begin{document}
\begin{center}
{\Large Dense sets with no nontrivial sum-product identity\par}
\medskip
{\large Counterexample to Conjecture 00000000038\par}
\end{center}

\begin{abstract}
For every integer $N\ge4$, the set $A_N=\{3,4,\ldots,N\}$ has density at
least $1/2$ in $[N]$ and admits no identity
$x_1+\cdots+x_k=x_1\cdots x_k$ of length $k\ge2$, even with repeated
choices. Thus Conjecture 00000000038 is false. The accompanying Lean 4
proof verifies the inequality for arbitrary finite lists, the exact
cardinality and density, arbitrarily large counterexamples, and the
negation of the conjecture even with an additional eventual-$N$ threshold.
\end{abstract}

\section{Statement and the nontriviality requirement}
The source asserts that for every $\delta>0$ there is a length $k(\delta)$
such that every $A\subseteq[N]$ with density at least $\delta$ contains a
nontrivial choice of $k$ elements, allowing repetition, whose sum equals
their product. Here $[N]=\{1,\ldots,N\}$ and density means $|A|/N$.

Both the English and Chinese source explicitly require nontriviality.
A one-term list always satisfies $x_1=x_1$, independently of $A$, so it
is the trivial case that must be excluded. We prove the stronger and
unambiguous fact that \emph{every} sum-product identity in the proposed
counterexample has length exactly one. Hence the counterexample applies
to any nontriviality convention excluding these singleton identities.
If singleton identities were admitted, taking $k=1$ would instead make
the assertion immediate and render the source's nontriviality clause
meaningless. Repetition is fully permitted in the proof below.

\section{An elementary obstruction}
\begin{lemma}\label{lem:inequality}
For every integer $k\ge2$ and every sequence of integers
$x_1,\ldots,x_k\ge3$,
\[
 \sum_{i=1}^k x_i < \prod_{i=1}^k x_i.
\]
\end{lemma}
\begin{proof}
If $u,v\ge3$, then
\[
 uv-u-v=(u-1)(v-1)-1\ge 2\cdot2-1=3>0.
\]
This proves the case $k=2$. Suppose now that $k\ge3$ and the result holds
for $k-1$ terms. Put
\[
 a=x_1,\qquad S=\sum_{i=2}^k x_i,\qquad
 P=\prod_{i=2}^k x_i.
\]
Then $a\ge3$, $S\ge6$, and $S<P$. Applying the two-variable inequality
to $a,S$ gives
\[
 a+S<aS<aP,
\]
which completes the induction. No distinctness of the $x_i$ was used.
\end{proof}

\begin{theorem}\label{thm:counterexample}
For every $N\ge4$, let $A_N=\{3,4,\ldots,N\}$. Then
\[
 A_N\subseteq[N],\qquad |A_N|=N-2,\qquad
 \frac{|A_N|}{N}=1-\frac2N\ge\frac12.
\]
For every $k\ge2$ and every choice $x_1,\ldots,x_k\in A_N$, with repetitions
allowed, the sum differs from the product.
\end{theorem}
\begin{proof}
The inclusion and cardinality follow directly from the endpoints of the
integer interval. Since $N\ge4$, $N-2\ge N/2$, proving the density claim.
Every selected element is at least $3$, so Lemma~\ref{lem:inequality}
gives a strict inequality and in particular excludes equality.
\end{proof}

Fix $\delta=1/2$. Any proposed nontrivial length $k(1/2)$ is excluded by
Theorem~\ref{thm:counterexample}. Moreover, the counterexamples are
arbitrarily large: for any proposed bound $B$, choose $N\ge\max\{B,4\}$.
Thus inserting a ``for all sufficiently large $N$'' qualification would
not repair the conjecture. Nor would allowing the chosen length to depend
on $A$ or $N$, because $A_N$ excludes all lengths $k\ge2$ simultaneously.

\section{Formal verification and scope}
The project pins Lean \texttt{v4.19.0} and Mathlib commit
\par\noindent
\texttt{c44e0c8ee63ca166450922a373c7409c5d26b00b}.
\par
The file \texttt{lean/Main.lean} contains the following kernel-checked
statements:
\begin{itemize}
\item \texttt{sum\_lt\_prod}\par
      The strict inequality for an arbitrary finite list with all entries
      at least $3$ and length at least $2$.
\item \texttt{only\_singleton\_solutions}\par
      Any such list whose sum equals its product has length exactly one.
\item \texttt{counterexampleSet\_half\_dense}\par
      The real-valued density inequality, proved from the separately
      checked exact interval cardinality.
\item \texttt{arbitrarily\_large\_dense\_counterexamples}\par
      A counterexample above every bound $B$ on the interval size.
\item \texttt{no\_uniform\_length\_even\_eventually}\par
      The negation of the fixed-density assertion, even after allowing a
      threshold $N_0$.
\end{itemize}

In the formalization, $A_N$ is the standard finite integer interval
\texttt{Finset.Icc 3 N}, and density is the actual real quotient
\texttt{(A.card : Real) / (N : Real)}. Selections are ordinary finite
lists of natural numbers, so repeated entries are retained rather than
removed. The final theorem has the logical shape
\[
 \neg\exists k\ge2\;\exists N_0\;\forall N\ge N_0,\ N>0\Longrightarrow
 \left[
 \begin{gathered}
  \text{every }A\subseteq[N]\text{ with }|A|/N\ge1/2\\
  \text{has a length-}k\text{ sum-product identity}
 \end{gathered}\right].
\]
This is a negation of a weaker assertion than the source's universal
claim, so it suffices to refute that claim. The only semantic point is
the source's explicit exclusion of trivial singleton identities; no
replacement arithmetic, density, or solution relation is used.

The full infinite family is proved by induction and exact algebra.
There is no finite search or numerical experiment standing in for a
universal proof. The Lean file contains no unfinished proofs or added
axioms. The printed axiom audit lists only
\texttt{propext}, \texttt{Classical.choice}, and \texttt{Quot.sound}.

\paragraph{Source.}
\href{https://github.com/The-Last-Math-Competition/The-Last-Math-Competition/blob/main/conjectures/00000000038.md}{The Last Math Competition, Conjecture 00000000038}.
The exact bilingual statement is preserved byte for byte in
\texttt{SOURCE.md}. Build instructions are in \texttt{README.md}; source
provenance and self-review are in the \texttt{verification} directory.

\end{document}
